\documentclass[11pt,a4paper]{article}

\usepackage[utf8]{inputenc}

\usepackage[T1]{fontenc}

\usepackage[polish]{babel}

\usepackage{lmodern}

\usepackage{geometry}

\geometry{margin=27mm}

\usepackage[hidelinks]{hyperref}

\title{Redukcja obciążenia poznawczego dzięki kontekstowej adaptacji interfejsu Smart Home}
% Tytuł katalogowy zachowano. Rodzaj planowanej publikacji:
% artykuł empiryczny z eksperymentem porównawczym (na razie konspekt).
% Autorstwo artykułu, kolejność autorów, afiliacje, korespondencja
% i identyfikatory ORCID wymagają potwierdzenia w projekcie kanonicznym.
\author{}

\date{Szkic naukowy, wersja 0.5.0}

\begin{document}
\maketitle

\begin{abstract}
Planowany artykuł empiryczny dotyczy związku między kontekstową adaptacją interfejsu Smart Home a obciążeniem poznawczym użytkowników. Celem proponowanego badania jest ustalenie, czy interfejs adaptacyjny różni się od statycznego pod względem wybranej miary obciążenia, a także czy znaczenie mają złożoność zadania, trafność rozpoznania kontekstu i przewidywalność zmian w interfejsie. Rozważany jest eksperyment z powtarzanymi pomiarami w symulacji Smart Home 2D, z równoważeniem kolejności warunków i zadań. Potencjalne zmienne obejmują subiektywne obciążenie, błędy wykonania oraz pomocnicze miary sprawności interakcji; hierarchia punktów końcowych i plan analizy wymagają zatwierdzenia. Koncepcja nie ma jeszcze potwierdzonego kanonicznego właściciela manuskryptu ani zamkniętego protokołu. Nie przeprowadzono badania i nie są dostępne wyniki empiryczne. Streszczenie opisuje projektowaną procedurę i pytania badawcze, a nie skuteczność proponowanego rozwiązania.
\end{abstract}

\noindent\textbf{Słowa kluczowe:} obciążenie poznawcze; adaptacja kontekstowa; inteligentny dom; interakcja człowiek--komputer; interfejs adaptacyjny; użyteczność; eksperyment z powtarzanymi pomiarami.

\medskip
\noindent\textbf{Rodzaj opracowania:} konspekt planowanego artykułu empirycznego; etap \texttt{IDEA} / \texttt{SOURCE\_PLANNED\_ARTICLE}.

\section{Wprowadzenie}
Systemy inteligentnego domu oferują wiele funkcji sterowania urządzeniami, scenami i powiadomieniami. Ich użytkownik musi rozpoznać bieżący cel, odnaleźć odpowiednią funkcję, zweryfikować stan urządzenia i ocenić skutki działania. Interfejs kontekstowy może ograniczać liczbę informacji wymagających przeszukiwania, ponieważ eksponuje funkcje związane z aktualnym zadaniem. Ta sama adaptacja może jednak zmienić położenie elementów, ukryć potrzebną opcję lub zaproponować działanie niedopasowane do intencji. Wtedy potencjalna oszczędność wysiłku może zostać zniwelowana przez konieczność ponownego zorientowania się w interfejsie.

Problem naukowy dotyczy zatem nie tylko różnicy między układem adaptacyjnym a statycznym, lecz także warunków, w których dopasowanie do kontekstu rzeczywiście pomaga użytkownikowi. Rozróżnić należy subiektywne obciążenie zadaniem, poprawność wykonania i szybkość obsługi. Zmiana któregokolwiek z tych wskaźników nie dowodzi automatycznie zmiany pozostałych. Proponowane badanie ma dotyczyć interakcji w symulacji Smart Home 2D, bez domniemania skuteczności w rzeczywistym gospodarstwie domowym.

\section{Podstawy teoretyczne i stan badań}
\subsection{Wymagania zadania a wysiłek związany z interfejsem}
Na potrzeby modelu pojęciowego odróżnia się wymagania wynikające z celu i liczby decyzji w zadaniu od dodatkowego wysiłku wywołanego sposobem prezentacji informacji. Przy podobnych zadaniach ograniczenie zbędnego przeszukiwania ekranu może zmniejszać wymagania obsługi, ale nie musi zmieniać złożoności samej czynności. Subiektywna ocena obciążenia zależy ponadto od oczekiwań, doświadczenia i przyjętej strategii działania. Błędy, czas wykonania i liczba operacji należy analizować jako odrębne miary, a nie zamienniki wysiłku umysłowego.

\subsection{Mechanizm selekcji informacji i mechanizm przeorientowania}
Rozważany mechanizm korzystny polega na wyeksponowaniu opcji przydatnych w danym kontekście i ograniczeniu konkurujących bodźców. Jego przewidywaną konsekwencją jest mniejsza potrzeba wyszukiwania i porównywania możliwości. Mechanizm przeciwny wiąże się z niestabilnością rozmieszczenia funkcji. Jeśli po zmianie kontekstu użytkownik musi ponownie szukać wcześniej poznanych elementów, adaptacja może zwiększać koszt przeorientowania i obniżać przewidywalność działania. Obydwa mechanizmy mogą działać równocześnie; ich względne znaczenie zależy prawdopodobnie od właściwości zadania i projektu interfejsu.

\subsection{Alternatywne wyjaśnienia potencjalnej różnicy między warunkami}
Nawet jeśli w przyszłym eksperymencie wystąpi różnica między warunkami, nie będzie ona sama w sobie dowodem zmniejszenia obciążenia poznawczego przez trafną selekcję informacji. Wynik może zależeć od różnej liczby koniecznych działań, wielkości i kontrastu przycisków, dodatkowych wskazówek wizualnych, znajomości wcześniejszego układu, kolejności prezentacji lub uczenia się obsługi. Dodatkowo błędne rozpoznanie kontekstu może obniżać skuteczność działania niezależnie od samego mechanizmu adaptacji. Projekt porównawczy musi zatem oddzielać jakość rozpoznania kontekstu od stabilności prezentacji i w miarę możliwości utrzymywać równoważność zadań oraz dostępnych funkcji.

\subsection{Luka badawcza i warunki nowości}
Samo przeciwstawienie interfejsu adaptacyjnego statycznemu nie stanowi wystarczającej przesłanki nowości. Właściwe pytanie dotyczy warunków brzegowych: kiedy trafne dostosowanie skraca wyszukiwanie informacji, a kiedy jego koszt, zwłaszcza przeorientowanie lub błąd kontekstu, może przewyższać korzyść. Proponowany wkład polegałby na rozdzieleniu skutków adaptacji od zmienności struktury i poprawności rozpoznania kontekstu, pod warunkiem że manipulacje te zostaną rzeczywiście rozdzielone w protokole. Zakres rzeczywistej luki i pierwszeństwo podobnych rozwiązań wymagają systematycznej weryfikacji literatury przed sformułowaniem twierdzeń o nowości.

\paragraph{Literatura wymagająca weryfikacji.}
Przed finalizacją podstaw teoretycznych należy sprawdzić publikacje dotyczące obciążenia poznawczego i obciążenia zadaniem w HCI, interfejsów adaptacyjnych Smart Home, stałości rozmieszczenia elementów, błędów adaptacji oraz zróżnicowania potrzeb użytkowników. Należy zweryfikować pełne teksty, miary, warunki eksperymentalne i granice uogólniania; dopiero wtedy przypisać konkretne cytowania do twierdzeń. W tej wersji nie przypisuje się niezweryfikowanych cytowań.

\section{Cel, pytania badawcze i hipotezy}
Celem proponowanego badania jest oszacowanie różnicy w uprzednio wybranej mierze obciążenia poznawczego między adaptacyjnym i statycznym interfejsem Smart Home przy porównywalnych zadaniach, a następnie zbadanie, w jakim stopniu różnica ta zależy od złożoności zadania, poprawności kontekstu i stabilności układu.

\begin{enumerate}
\item \textbf{PB1:} Czy podczas wykonywania porównywalnych zadań interfejs adaptacyjny różni się od statycznego pod względem wybranego wskaźnika obciążenia poznawczego?
\item \textbf{PB2:} Czy złożoność zadania zmienia różnicę między warunkami?
\item \textbf{PB3:} Czy nieadekwatna adaptacja lub przebudowa układu wiąże się z dodatkowymi kosztami przeorientowania i wykonania zadania?
\end{enumerate}

\subsection{Kandydacka hipoteza konfirmacyjna}
Jeżeli w protokole zostanie zapewnione poprawne rozpoznanie kontekstu, porównywalny zakres funkcji i stabilne rozmieszczenie istotnych elementów, można rozważyć hipotezę \textbf{H1}, zgodnie z którą adaptacja selekcji informacji obniża wcześniej określony wskaźnik subiektywnego obciążenia w stosunku do interfejsu statycznego. Jest to \emph{hipoteza projektowa}, a nie zatwierdzona hipoteza konfirmacyjna: wymaga niezależnego uzasadnienia w zweryfikowanej literaturze, wyboru jednej miary pierwszorzędowej, jednoznacznego kontrastu i zamknięcia protokołu przed analizą danych. Jeśli te warunki nie zostaną spełnione, porównanie należy traktować jako eksploracyjne, bez kierunkowej deklaracji.

\subsection{Pytania i przewidywania eksploracyjne}
Dla \textbf{PB2} planuje się eksploracyjne sprawdzenie interakcji między rodzajem interfejsu a złożonością zadania. Dla \textbf{PB3} można ocenić, czy błędne rozpoznanie kontekstu lub zmiana lokalizacji funkcji zwiększa liczbę pomyłek, czas ponownego odnalezienia opcji albo subiektywne obciążenie. Kierunek i wielkość tych zależności nie są z góry przesądzone. Jeżeli protokół nie będzie manipulował osobno trafnością kontekstu i stabilnością układu, ich efekty nie powinny być interpretowane przyczynowo.

\subsection{Warunki rozstrzygnięcia i zakres wnioskowania}
Przed rekrutacją należy określić pierwszorzędowy punkt końcowy, estymand, znaczenie praktyczne różnicy, jednostkę analizy, reguły wykluczeń i plan kontroli wielokrotności. Dopasowanie scenariuszy powinno zapobiegać myleniu adaptacji z redukcją liczby kroków lub inną zmianą ergonomii. Przewidywania nie stanowią prerejestracji, a wyniki w symulacji 2D nie uzasadnią automatycznie wniosków dotyczących osób starszych, ograniczeń poznawczych ani funkcjonowania w rzeczywistym domu.

\section{Projekt badania i metody}
\subsection{Status protokołu i charakter stanowiska}
Poniższa procedura jest projektem badania, nie opisem wykonanych czynności. Zakłada interakcję człowieka z symulacją inteligentnego domu w dwuwymiarowym interfejsie ekranowym. Do podstawowego porównania nie są wymagane EEG, okulografia ani XR. Zastosowanie takich pomiarów wymagałoby osobnej hipotezy, procedury, walidacji aparatury oraz oceny etycznej; nie stanowią one części obecnego protokołu. Planowane stanowisko \texttt{ST-06} nie zostało zweryfikowane, a aplikacja \texttt{app-carehome} ma status projektu niewdrożonego. Nie zakłada się dostępności kompletnego oprogramowania ani aparatury.

\subsection{Schemat eksperymentu i warunki}
Rozważany jest wewnątrzosobniczy eksperyment porównujący interfejs statyczny z adaptacyjnym w symulacji Smart Home 2D. W obu warunkach należy zachować jednakową pulę funkcji, cele zadań, komunikaty o rezultacie i zasady nawigacji. Warunek statyczny utrzymywałby niezmieniony sposób prezentacji opcji; warunek adaptacyjny eksponowałby funkcje wskazane przez jawne reguły kontekstowe, bez zmiany wymagań zadania. Aby przetestować H1, należy ograniczyć przeorganizowanie położenia elementów, gdyż w przeciwnym wypadku efekt selekcji informacji byłby nierozdzielny od kosztu przeorientowania.

Poza porównaniem podstawowym rozważa się zróżnicowanie złożoności scenariuszy oraz eksploracyjne warunki błędnego rozpoznania kontekstu i zmiennego układu. Takie czynniki można interpretować przyczynowo tylko po ich niezależnej manipulacji i przypisaniu prób do warunków; przy samym pomiarze obserwacyjnym pozostają analizami asocjacyjnymi. Liczba poziomów i kolejność manipulacji wymagają ustalenia przed pilotażem.

\subsection{Populacja i dobór próby}
Populacja docelowa pozostaje do rozstrzygnięcia. Wstępnie można rozważać osoby pełnoletnie zdolne do świadomej zgody i samodzielnej obsługi badanego urządzenia wejściowego. Proponowane kryteria włączenia obejmują możliwość przeczytania instrukcji w języku badania i wykonania zadań próbnych, również z uzgodnionymi udogodnieniami dostępności. Potencjalne kryteria wyłączenia powinny dotyczyć wyłącznie ograniczeń uniemożliwiających wiarygodny pomiar mimo udogodnień, braku zgody lub niespełnienia jawnych kryteriów protokołu. Nie należy wykluczać osób ze względu na wiek albo niepełnosprawność bez uzasadnienia zgodnego z celem badania.

Należy udokumentować sposób rekrutacji, doświadczenie z urządzeniami Smart Home, podstawową sprawność obsługi oraz istotne czynniki mogące różnicować trudność zadań. Wielkość próby wymaga oszacowania na podstawie wybranego estymandu, minimalnego efektu istotnego praktycznie, wariancji i korelacji powtarzanych obserwacji albo założonej precyzji przedziału ufności. W tym szkicu nie deklaruje się liczby uczestników, przeprowadzonej rekrutacji ani oszacowanej mocy.

\subsection{Konstrukty i operacjonalizacja}
Głównym konstruktem jest obciążenie poznawcze związane z obsługą interfejsu przy wykonywaniu zadania. Kandydatem na jego wskaźnik jest subiektywna ocena wysiłku umysłowego po określonej jednostce pracy, na przykład odpowiedni wymiar zweryfikowanego narzędzia do pomiaru obciążenia zadaniem. Przed badaniem należy wybrać konkretny instrument, wersję językową, zasady punktowania, moment podania i uzasadnić jego trafność w tej populacji. Nie należy traktować pełnego wyniku skali wielowymiarowej automatycznie jako czystej miary obciążenia poznawczego.

Zmienną niezależną w porównaniu podstawowym będzie rodzaj interfejsu (adaptacyjny lub statyczny). Złożoność zadania wymaga obiektywnej definicji opartej na wymaganych decyzjach, zależnościach czynności i liczbie informacji potrzebnych do osiągnięcia celu, a nie wyłącznie na czasie realizacji. Trafność kontekstu należy określać jako zgodność klasyfikacji sytuacji z wcześniej oznaczonym stanem referencyjnym. Stabilność układu można mierzyć liczbą i zakresem przemieszczeń tych samych funkcji między widokami.

Wskaźniki pomocnicze obejmowałyby wykonanie celu (tak/nie), liczbę błędów według zamkniętego katalogu, czas od rozpoczęcia próby do potwierdzenia osiągnięcia celu oraz liczbę operacji nawigacyjnych. Błąd krytyczny powinien oznaczać wcześniej zdefiniowane niepożądane działanie w symulacji, a nie ocenę dokonaną po poznaniu warunku eksperymentalnego. Wyniki behawioralne nie zastępują pomiaru obciążenia. Przed akwizycją należy wybrać jeden pierwszorzędowy punkt końcowy albo z góry uzasadnić wyniki współpierwszorzędowe i sposób kontroli wielokrotności.

\subsection{Scenariusze i procedura uczestnika}
Przykładowe klasy zadań do przyszłego opracowania obejmują ustawienie parametrów urządzenia, uruchomienie sceny i sprawdzenie stanu systemu. Każdy scenariusz powinien mieć jednoznaczny cel, warunki początkowe, możliwe ścieżki wykonania, reguły zaliczenia i poziom trudności. Zestawy należy porównać pod względem równoważności, aby różnica wyników nie wynikała jedynie z innej długości lub dostępności ścieżki.

Po uzyskaniu wymaganych zezwoleń i zgody uczestnik zapoznałby się z instrukcją oraz przeszedł neutralny trening obsługi. Następnie wykonywałby zadania w obu warunkach, z kontrbalansowaniem kolejności interfejsów i zestawów zadań oraz z ujednoliconymi przerwami. Po zdefiniowanym bloku albo zadaniu oceniałby obciążenie wybranym narzędziem. Dziennik eksperymentu rejestrowałby identyfikator pseudonimowy, wersję symulacji, warunek, scenariusz, znacznik czasu, akcje i rezultat bez zbędnych danych osobowych. Dokładne momenty pomiaru oraz zasady przerywania próby muszą zostać zapisane przed rekrutacją.

\subsection{Pilotaż i kontrola jakości}
Przed badaniem właściwym konieczny jest pilotaż zadań, czytelności instrukcji, równoważności obu interfejsów, działania reguł adaptacji i kompletności rejestrowanych zdarzeń. Testy techniczne powinny obejmować powtarzalność stanów początkowych, zgodność logów z działaniami, poprawność znaczników czasu, brak ujawniania warunku przez instrukcję oraz kontrolę awarii i przerwanych prób. Reguły odrzucania prób z powodu błędu technicznego, nieukończenia zadania, brakującej oceny i naruszenia procedury należy ustalić przed analizą, z rozróżnieniem braków technicznych i zachowania uczestnika.

Scenariusze, definicje zmiennych, reguły adaptacji, konfiguracje obu interfejsów, wersje oprogramowania oraz protokół i plan analiz powinny zostać utrwalone przez kanoniczny projekt badawczy. Nie przeprowadzono tu pilotażu ani walidacji pomiaru.

\subsection{Etyka, bezpieczeństwo i przetwarzanie danych}
Przed rozpoczęciem rekrutacji należy ustalić wymagania właściwej komisji etycznej i uzyskać odpowiednie rozstrzygnięcie, jeśli jest wymagane. Uczestnicy powinni otrzymać zrozumiałą informację o badaniu, dobrowolności udziału, prawie wycofania się, danych gromadzonych w logach i zasadach ich przechowywania. Projekt powinien uwzględnić dostępność interfejsu, zmęczenie i przerwy, a działania w symulacji nie mogą sterować rzeczywistymi urządzeniami bez odrębnej procedury bezpieczeństwa.

Należy stosować minimalizację danych, pseudonimizację identyfikatorów, oddzielne przechowywanie ewentualnego klucza identyfikacyjnego, kontrolę dostępu i ustalony czas retencji. Nie deklaruje się otrzymania zgody etycznej, pozyskania danych ani udostępnienia zbiorów. Rozstrzygnięcie właściciela kanonicznego manuskryptu i protokołu pozostaje warunkiem dalszych prac.

\section{Plan analizy danych}
\subsection{Estymand i porównanie podstawowe}
Przed akwizycją danych należy wyznaczyć kontrast adaptacja--interfejs statyczny dla wybranego wyniku pierwszorzędowego, jednostkę analizy i sposób uwzględnienia zadań powtarzanych. Kandydatem na metodę jest model mieszany obejmujący efekty uczestnika i zadania. Rodzina rozkładu i funkcja łącząca będą zależeć od własności skali wynikowej; nie należy automatycznie stosować modelu liniowego do liczby błędów.

\subsection{Moderacja i czynniki zakłócające}
Plan powinien obejmować co najmniej ocenę interakcji warunku z poziomem złożoności zadania, kontrolę kolejności ekspozycji oraz analizę efektu uczenia. Ocenę skutków nieadekwatnej adaptacji trzeba uprzednio powiązać z manipulacją lub jasno określić jako analizę eksploracyjną.

\subsection{Kontrola jakości i analizy wrażliwości}
Należy z góry określić kryteria poprawności prób, traktowanie braków danych, reguły wykluczeń, diagnostykę dopasowania modelu i zakres analiz wrażliwości. Raportowanie powinno uwzględniać oszacowania efektów, przedziały niepewności i ograniczenia interpretacji. Nie wykonano obliczeń ani nie oszacowano parametrów empirycznych.

\section{Zakres przyszłego raportowania wyników}
Po zebraniu i zweryfikowaniu danych część wynikowa powinna przedstawić charakterystykę próby, kompletność obserwacji, jakość manipulacji, porównanie pierwszorzędowe oraz ustalone analizy uzupełniające. Na obecnym etapie nie istnieją wyniki, które można byłoby w tej sekcji przedstawić. Rozdział zawierający rezultaty empiryczne zostanie dodany wyłącznie po ich uzyskaniu.

\section{Ramy przyszłej dyskusji}
Interpretacja będzie wymagała oddzielenia redukcji subiektywnego obciążenia od poprawy szybkości i trafności działania. Należy rozważyć, czy ewentualna korzyść wynika z selekcji informacji, czy z innych różnic między wersjami interfejsu, oraz czy koszty przeorientowania równoważą korzyści adaptacji. Dyskusja nie może przedstawiać planowanych zależności jako potwierdzonych.

\section{Ograniczenia koncepcji}
Najważniejsze ograniczenia planu obejmują możliwą nietrafność operacjonalizacji obciążenia, nierównoważność zadań, efekty uczenia, błędy rozpoznania kontekstu i zależność wyników od doświadczenia użytkowników. Badanie w środowisku 2D nie będzie samo w sobie dowodem skuteczności w rzeczywistym Smart Home ani poprawy codziennej samodzielności. Własność kanonicznego manuskryptu oraz kompletność dokumentacji źródłowej pozostają niepotwierdzone.

\section{Wnioski na etapie planowania}
Zaproponowano strukturę przyszłego badania porównawczego, która rozdziela obciążenie poznawcze, sprawność wykonania i potencjalne koszty adaptacji. Przed rozpoczęciem realizacji konieczne są przegląd literatury, ustalenie kanonicznego właściciela projektu, wybór głównego punktu końcowego, zamknięcie protokołu i potwierdzenie warunków etycznych. Nie formułuje się wniosków o efektywności interfejsu.

\section*{Metadane autorstwa i oświadczenia wydawnicze}
Autorstwo i kolejność autorów, afiliacje instytucjonalne, adres autora korespondencyjnego oraz identyfikatory ORCID pozostają do weryfikacji w repozytorium kanonicznym. Puste pole \verb|\author{}| nie oznacza pracy anonimowej ani rozstrzygnięcia autorstwa.

Deklaracje wkładów według taksonomii CRediT należy przygotować po potwierdzeniu osób i ich rzeczywistego udziału w konceptualizacji, metodologii, oprogramowaniu, analizie, przygotowaniu tekstu oraz nadzorze, odpowiednio do faktycznego wkładu. Nie przypisano obecnie żadnej roli konkretnej osobie.

Informacje o źródłach finansowania i ewentualnej roli podmiotów finansujących wymagają potwierdzenia. Brak tych informacji w szkicu nie stanowi deklaracji braku finansowania.

Przed zgłoszeniem artykułu należy odrębnie zweryfikować wymagane przez wybrane czasopismo oświadczenia dotyczące konfliktów interesów, oceny etycznej i zgody uczestników (jeżeli dotyczą badania), dostępności danych i kodu, korzystania z narzędzi AI, praw do materiałów oraz zgodności z zasadami wydawniczymi. Żadnego zatwierdzenia ani udostępnienia materiałów nie potwierdzono.

\section*{Status koncepcji i źródeł}
Praca ma status \texttt{IDEA} i \texttt{SOURCE\_PLANNED\_ARTICLE}. Katalog ARC Neuro pełni funkcję indeksu, a rejestry \texttt{portfolio/projects.json} i \texttt{portfolio/papers.json} nie wskazują obecnie kanonicznego właściciela manuskryptu. Literatura wymaga weryfikacji i przypisania cytowań przed przygotowaniem publikacji. Brak zatwierdzonego protokołu, danych i wyników nie oznacza wyniku zerowego.

\end{document}
